% SGA 3，第 VII_B 篇，英文重建稿——第 1 轮。
% 范围：权威底本分篇第 91 页，第 5.2.1.B 节全文，止于第 5.2.1.C 节之前。
% 合并阅读本页码： 604.
% 再版印刷页码： 593.
% 本范围内编者注： 156.
% 权威底本： Polo--Gille Exp7B-13oct24.pdf,
% SHA-256 D303690E0195654BBE7C3DAA5D89EA23E24CD709A3575217234AC19AE8836694.
% 曾参考 Jacob Reinhold 的 CC BY 4.0 Markdown 作为比对谱系。
% 本组件没有图表；全部数学均为原生 LaTeX。

\paragraph{\SGARefLink{sga3:VIIB:5.2.1.B}{5.2.1.B}.}
\phantomsection
\label{sga3:VIIB:5.2.1.B}

如上，置
\[
  \omega_\infty
  =
  \overline{\bigcup_{n\in\mathbb N}\omega_n}.
\]
所需类型的截面
\[
  \sigma:\omega\longrightarrow I
\]
存在的第二种情形，是 \(\omega_\infty\) 有一组伪基 \(\mathscr B_\infty\)，它是各空间
\[
  \omega_n/\omega_{n-1}.
\]
的伪基在 \(n\in\mathbb N^*\) 上的并。此时可用 \(\omega/\omega_\infty\) 的伪基 \(\mathscr B'_\infty\) 将它补成 \(\omega\) 的一组与滤过相容的伪基。

置
\[
  V=\omega_\infty^\dagger,
\]
令 \(V_n\) 为 \(\omega_n\) 在 \(V\) 中的正交补。上述条件等价于在通常的 \(k\)-向量空间范畴中，分离递减滤过
\[
  V=V_0\supset V_1\supset V_2\supset\cdots
\]
分裂：即 \(V\) 为子空间 \(F_n\) 在 \(n\in\mathbb N\) 上的直和，且满足
\[
  F_n\simeq V_n/V_{n+1}.
\]

这未必成立。例如，令
\[
  V=\mathscr S=k^{\mathbb N}
\]
为 \(k\) 中元素的序列空间，令 \(\mathscr S_n\) 为满足当 \(i<n\) 时 \(u_i=0\) 的序列 \((u_i)\) 所组成的子空间。则
\[
  \dim\mathscr S_n/\mathscr S_{n+1}=1,
\]
但 \(\mathscr S\) 不同构于各 \(\mathscr S_n/\mathscr S_{n+1}\) 的直和，因为 \(\mathscr S\) 并非可数维的。不过，在此例中，\(\mathscr S\) 是各 \(\mathscr S_n/\mathscr S_{n+1}\) 的积\textup{（参见 \SGARefLink{sga3:VIIB:5.2.1.A}{5.2.1.A}）}。

当 \(V\) 为可数维时，分裂确实成立。\footnote{%
\label{sga3:VIIB:nde:156}此段来自与 J.-M. Fontaine 和 E. Bouscaren 的讨论；特别是 Bouscaren 指出了下面的证明。}
事实上，设
\[
  (e_n)_{n\in\mathbb N}
\]
为 \(V\) 的一组基。我们将归纳构造递增函数
\[
  g:\mathbb N\longrightarrow\mathbb N
\]
以及对 \(i=0,\ldots,g(n)\) 的子空间 \(F_i\)，使得
\[
  F_i\simeq V_i/V_{i+1}
\]
且
\[
  F_{\leq g(n)}
  =
  \bigoplus_{i=0}^{g(n)}F_i
\]
为 \(V_{g(n)+1}\) 的一个补空间，并包含
\[
  e_0,\ldots,e_n.
\]
于是可得
\[
  V=\bigoplus_{i\geq0}F_i.
\]

设 \(n+1\in\mathbb N\)，且已对 \(n\) 证明了断言\textup{（当 \(n=-1\) 时为空断言）}。若
\[
  e_{n+1}\in F_{\leq g(n)},
\]
置 \(g(n+1)=g(n)\)。否则写成
\[
  e_{n+1}=f+x,
  \qquad
  f\in F_{\leq g(n)},
  \qquad
  0\ne x\in V_{g(n)+1}.
\]
令 \(j\) 为使得
\[
  x\in V_j\setminus V_{j+1}.
\]
成立的最小整数。对
\[
  i=g(n)+1,\ldots,j,
\]
选取 \(V_{i+1}\) 在 \(V_i\) 中的补空间 \(F_i\)，并使所选 \(F_j\) 满足 \(x\in F_j\)。则
\[
  e_{n+1}\in F_{\leq j};
\]
置 \(g(n+1)=j\)。

\phantomsection
\label{sga3:VIIB:5.2.1.B:induction-source-note}
\noindent\textit{底本说明。}\enspace
在归纳的最后一步，法文 PDF 印作 \(e_{n+1}\in F_j\)。由 \(e_{n+1}=f+x\) 及 \(f\in F_{\leq g(n)}\)，论证实际给出 \(x\in F_j\) 和 \(e_{n+1}\in F_{\leq j}\)；这里采用这些关系。
